% !TEX root = main.tex
\section{보안 증명}
정확성, 건전성, 프라이버시를 차례로 증명한다. 원장의 권한·
등록 판정과 원자적 갱신, 불변 생성·소비 기록, 확정 시점의 인
증 조회, 신뢰된 감사자와 상태 관리자의 집행은 모델 전제다.
모든 소비 경로는 실행 시점의 미소비·활성 상태를 검사한다.
키의 외부 전달은 유효한 AuditKey를 제공하는 신뢰 경계로 두

% Source PDF page 12, column 1
며, 그 서비스의 보안을 새로 증명하지 않는다. 암호 계약과 확
률은 전체 다항 실행 예산에 대한 것이고 실제 방식의 가정 충족
여부는 별도 검증 대상이다.

\subsection{정확성(Correctness)}
정확성 실험은 정직한 생성 증인을 보존한다. 전이 호출에는 전
체 관계·실행 조건·신규 출력·저장 용량이 유효하고 집행 전까
지 유지되는 상태를, 공개 DPP 검증에는 발급이 포함된 인증 이
력과 일치하는 문서·주장을 요구한다. 추적에는 유효 시작점·고
정 \(H_{t}\)·유효 키와 필요한 유한 기록, 종료하는 인증 조회를 요구
한다. 정상 조건이 성립하는 호출에서 기대 결과를 내지 못하면
Failcorr로 기록한다.

동결의 정상 조건은 별도로 둔다. 차단 승인이 있고 정확한
전체 대상 \(T\) 가 비어 있지 않으면, \(H_{t}\)부터 공통 \(H_{r}\)까지 대상이
소비되지 않고 이미 차단된 상태는 유지되어야 한다. Active 대
상의 합법적 동결, 모든 제출 결과와 조회, 사용한 관찰을 반영
하는 같은 확정 이력의 \(H_{r}\) 확인이 설정된 유한 한도 안에 끝나
며 차단 해제가 없어야 한다. 이 조건 없이 완료까지 약속하지
않는다.

정리 1 (Correctness). 정직 증명 생성이 \(Q\)회 이하이고 1회
완전성 실패 상계가 \(\epsilon _{comp}\)이면

\begin{equation}
    \Pr[\mathrm{Fail}_{\mathrm{corr}}]\le Q\epsilon_{\mathrm{comp}}+\epsilon_E+\epsilon_H.\tag{23}
\end{equation}
\(\epsilon _{E}\)는 전체 복호화 실패, \(\epsilon _{H}\)는 관련 해시 충돌의 전체 실행 상계
로 다음 건전성 절의 계약을 사용한다. 여기에는 정직 생성에서
도 가능한 서로 다른 객체의 typed 소비 식별자 충돌이 포함된
다. 출력 참조의 신규성만으로 그 충돌까지 배제하지 않는다.
정상 조건과 위반이 함께 발생하는 사건의 상계이며, 정상 조건
에 조건을 건 확률 상계와 구분한다.

증명. 완전성 실패 밖에서 정직 증명은 전체 관계를 만족한다.
나머지 실행 조건도 유효하므로 전이는 수락되고 결정적 Apply
가 상태와 색인을 갱신한다. 동일 정규 인코딩과 등록 Claim을
쓰는 공개 DPP도 valid로 판정된다. 정직 생성에서 감사 암호
문은 실제 부모·미래 식별자를 담으므로 복호화 실패 밖에서 그
값을 얻는다. 관련 충돌이 없으면 객체의 소비 식별자는 구별된
다.

따라서 9.2절의 경로 귀납을 정직한 생성 목록에 적용하여
정확한 선행·후속 그래프·경계·종단을 얻는다. 유한 그래프·방
문 처리와 종료하는 조회로 탐색도 끝난다. 고정 \(H_{t}\)의 추적에
는 그 이후 원장 전체의 무경합을 요구하지 않는다. 정직 증인
을 이미 아는 이 실험에 지식 추출은 필요하지 않다.

동결은 전체 \(T\) 를 보존하고, \(T\) 가 비면 빈 대상 분기를 따른다.
그 밖에서는 위 처리·확정· 상태 유지 조건으로 같은 \(H_{r}\)에서 모
든 대상이 미소비·차단이며 대상별 결과와 제출도 정산된다. 따
라서 완료 술어가 참이 되어 COMPLETE\_AT\_CHECKPOINT를
반환한다. 남은 실패는 완전성· 복호화 실패와 관련 충돌의 합
집합에 포함되어 상계가 따른다. \(□\)

\subsection{건전성(Soundness)}
실제 전이를 기준으로 한 게임과 계약. 보안 매개변수 \(λ\)에 대
해 다항 길이 실행을 고려한다. 공격자는 정직한 전이 생성과
임의의 \((x,π)\) 제출을 적응적으로 섞으며, 임의의 확정 prefix와
문서·Claim을 공개 검증하는 CheckDPP도 질의한다. 원장은
고정된 전체 관계의 검증키로 증명을 확인하고 현재 권한·정책·

% Source PDF page 12, column 2
미소비·활성·기한·출력 신규성을 검사한 뒤 원자적으로 갱신한
다. 수락 이력을 처음부터 의미적으로 유효하다고 가정하지 않
는다.

분석용 생성 목록 \(C\)에 각 객체가 처음 수락될 때의 전체 개방
값과 typed 미래 식별자를 고정한다. 정직 생성에는 원래 증인,
공격자 수락에는 추출 증인을 사용하며 식별자에서 생성 객체
로 가는 역대응도 고정한다. 이는 구현에 추가하는 공개 목록이
아니다. 실패 사건 BadSem은 수락 입력이 먼저 생성되지 않았
거나, 생성 개방값과 다른 값을 사용하거나, 그 실제 값에 대한
소유 비밀· 정수 회계·정책 관계 또는 현재 상태 검사가 거짓인
경우다. 반복 계상·소비, 잘못된 상태 갱신, 실제 부모·미래 식
별자와 다른 감사 복호화도 포함한다. 공개 DPP는 인증된 발급
기록과 입력 문서 해시·정책·논스의 불일치를 실패로 정의한다.
추출 실패나 모호한 역대응에서는 분석 게임이 실패를 기록하
고 종료한다. 그 밖에서만 실제 생성 목록과 의미 위반을 비교
한다.

공격자는 수락 이력의 유효 시작점·확정 시점·방향을 골라
추적을 실행하게 할 수 있다. 차단 승인이 있는 후속 추적에는
신뢰된 상태 관리자가 AuditAndFreeze를 수행한다. 공격자는
허용된 전이· 합법적 상태 변경·자료 오류를 스케줄하며 감사
이후의 전이까지 전체 실행 예산에 포함한다. 인증 조회는 실제
값 또는 오류만 제공하고, 임의 권한 변경이나 새로운 키 공개
질의는 주지 않는다.

통합 게임 Expsound는 위 전이·DPP 의미 위반, 잘못된 추적
성공, 또는 완료 반환과 완료 술어/차단 유지 구간의 위반이 함
께 생기면 1을 반환한다. 자료 부재의 명시적 실패나 전체 대상
을 보존한 INCOMPLETE는 공격 성공이 아니다. 아래 나쁜 사
건에서는 1로 종료하고 그 밖에서만 생성 목록과 결과를 비교한
다.

다음은 해당 관계를 구현하는 방식의 충분한 계약이다.

1. 정직한 proof 관측·재전송과 다수 등록 관계를 포함하

는 적응적 실행의 수락 증인을 공동 추출하는 PPT 알고

리즘이 있다. 비공개 추출 trapdoor 또는 명시된 공격자

접근 권한을 가지며 이를 관측자에게 주지 않는다. 공개

transcript 분포는 실제 실행과 같거나 통계적으로 가까우

며 그 오차도 전체 실패 상계 \(\epsilon _{X}\)에 포함한다. 추출 보조정

보는 독립 원시기술의 challenge 비밀을 주지 않아야 한

다. 정직한 동일 증명 재전송에는 원래 증인을 쓴다. 단일

증명의 지식 건전성만으로 이 계약을 추정하지 않는다.

2. 정규 전체 개방값의 계산적 결합성, accepted-root 포함

관계의 건전성, 관련 객체·식별자·Claim 해시의 충돌 저항

성에 대한 전체 실행 상계를 각각 \(\epsilon _{C},\epsilon _{M},\epsilon _{H}\)라 둔다. 유

효 비밀 영역에서 OwnerKey는 단사이다.

3. 고정 키·문맥에서 모든 허용 평문과 prover가선택한 유효

난수에 대해 복호화가 정확하다. 완전 정확성이면 실패는

0이며, 완화할 경우 적대적 난수 선택의 실패 상계를 \(\epsilon _{E}\)로

둔다. 정직 난수 평균에서만 정확하다는 조건은 충분하지

않다.

소유권 키의 예로 소수 위수 \(\ell \)의 생성점 \(G\)에 대해 \(sG\)와 정규 영
역 \(1≤s<\ell \)를 사용할 수 있다. 단사성과 비밀의 일방향성은
다른 성질이다. 원장·인증·권한·정책 구성의 정확성은 모델 전
제이며 위 암호 가정으로 새로 증명하지 않는다.

정리 2 (Soundness). 위 계약 아래

\begin{equation}
    \Pr[\mathrm{Exp}_{\mathrm{sound}}=1]\le\epsilon_B:=\epsilon_X+\epsilon_C+\epsilon_M+\epsilon_H+\epsilon_E.\tag{24}
\end{equation}

% Source PDF page 13, column 1
각 항은 전체 다항 실행에 대한 환원 상계다. 소유권 결론은 관
계에 맞는 비밀의 지식이며 그 비밀의 누출 방지나 현실 입력의
진실성까지 뜻하지 않는다.

증명: 전이와 감사 결합. 추출·결합·포함 관계·관련 충돌·복
호화 계약의 실패 합집합을 Bad라 두면 Pr\([\)Bad\(]≤\epsilon _{B}\)다. 그 밖
에서 수락 순서로 귀납한다. 추출 입력의 포함 증인은 이전 실
제 등록 객체를 가리킨다. 생성 목록의 개방값과 다르면 계산적
결합성의 명시적 충돌 쌍이다. 따라서 추출 증인의 타입·정수·
소유권·정책 관계가 바로 최초 생성 값에 적용된다. 계산적 결
합성을 모든 수학적 증인의 유일성으로 해석하는 것은 아니다.

Note의 생성·소비 개방값은 같은 \(cm,owner\)를 포함한다.
단사성으로 소유 비밀도 같아 \(nf\)가 같다. Voucher는 결합된
같은 개방 난수에서 \(rvnf\)를 얻으므로 인수·회수도 동일하다.
다른 같은 타입 객체의 식별자 중복은 서로 다른 참조가 포함
된 정규 해시 입력의 충돌이다. 이에 따라 현재 typed spent/
status 검사가 바로 실제 객체에 적용된다. 중복 입력 검사와
원자적 소비 기록은 반복 계상과 경합 재소비를 거부한다. 과거
membership root는 현재 상태 검사를 대체하지 않는다.

감사 관계는 이 실제 입력·출력의 부모·미래 식별자를 같은
증인으로 암호화한다. 고정 키·문맥·슬롯과 적대적 난수의 복호
화 정확성으로 감사자는 그 평문을 얻는다. 다른 문서 해시·정책
·논스를 기존 등록 Claim에 붙여 유효 판정을 받으면 Claim 결
합 해시의 충돌이다. 따라서 실패 합집합 밖에서는 BadSem이
없고, 같은 나쁜 사건 밖에서 추적·차단을 분석할 수 있다. 아래
두 보조논증으로 정리 2의 나머지 결과를 보인다.

실제 그래프와 추적 보조논증. 생성 목록과 typed 식별자 역
대응으로 실제 입력을 고정하고, 객체에서 실제 소비 이벤트로,
이벤트에서 전체 출력으로 간선을 잇는다. 생성이 소비보다 앞
서고 생성이 유일하므로 DAG다. 이를 감사 암호문이나 반환
결과로 정의하지 않는다. 앞의 전이 논증에서 제외한 나쁜 사
건 밖에서 TraceForward가 정상 반환하면, 그 결과는 \(H_{t}\)의 지
정 시작점에서 도달한 소비 이벤트의 모든 출력을 전개하는 후
속 그래프 및 그 미소비 회계 증표·미해결 이전 증표 집합과 일
치한다. 단일 객체 시작은 해당 객체에서 시작하고, 기록 시작
은 그 기록의 모든 출력에서 시작한다. 기록 시작의 후속 결과
에는 seed 기록 자체와 seed 출력 간선을 추가하지 않으므로
Exit 시작은 빈 결과다. 선행 기록 시작은 그 기록과 직접 출력
·출력 간선을 포함한 뒤 부모를 거슬러 올라간다. 재귀 중 읽
은 다른 producer의 sibling은 넣지 않는다. 이는 단일한 event
reachability 정의와 다르다. 합류한 소비 이벤트의 다른 부모
에 대한 선행 그래프까지 자동 포함하는 명제는 아니다.

이유는 시작점으로부터의 경로 길이에 대한 귀납으로 설명
할 수 있다. 실제 출력과 결합된 미래 식별자로 조회하므로, 소
비 기록이 0이면 정확히 해당 시점의 미소비 객체이다. 0이 아
니면 불변 소비 기록이 유일한 실제 소비 이벤트를 가리키며, 타
입과 공개 입력을 확인한 뒤 그 소비 이벤트의 모든 출력을 전
개한다. 이 과정을 반복하면 정의된 후속 경로를 빠뜨리지 않는
다. 방문 기록은 중복 전개를 줄이되 도달 간선은 보존해야 한
다. 제품 주장과 종료는 정의된 종단이며, 유한 그래프와 자료
가용성 아래에서 탐색이 끝난다. 선행 추적도 실제 부모의 결합
과 불변 생성 기록 조회에 대해 같은 방식으로 설명할 수 있다.
자료 누락 시의 실패는 이 정상 반환 성질과 양립하며 무조건적
인 가용성을 의미하지 않는다.

형식적으로 공격자가 수락 이력 뒤 확정 prefix·유효 시작점·
방향을 고르고 자료의 누락·오류를 유발하게 한다. 인증된 조회

% Source PDF page 13, column 2
는 그 prefix의 실제 값 또는 오류만 제공하고 감사자에게 유효
한 키를 준다. 성공 결과가 실제 생성 목록의 위 그래프·경계·종
단과 다르면 공격 성공이다. 앞의 귀납은 같은 나쁜 사건 밖에
서 이를 배제하므로 실패 확률은 \(\epsilon _{B}\) 이하다. 정상 종료는 자료
가용성에 별도로 조건부이며 나쁜 사건 밖에서 성립한다. 분석
게임은 나쁜 사건에서 실패를 기록하고, 그 밖에서만 목표 그래
프를 비교한다. 항상 실패하는 절차를 정확한 감사로 취급하지
않는다.

이 논증은 검사 생략의 영향을 드러낸다. 출력과 무관한 식
별자를 암호화할 수 있으면 추적은 다른 객체의 소비 여부를 조
회할 수 있고, 인증되지 않은 0 응답을 받아들이면 이미 소비된
객체를 미소비 경계 집합의 원소로 잘못 분류할 수 있다. 일부
소비 경로가 상태 검사를 생략하면 정확한 추적과 동결만으로
그 경로의 소비를 막을 수 없다. 이는 각 검사의 역할을 설명하
는 모델 반례이며 구체 구현에서 그러한 취약점이 발견되었다
는 보고는 아니다.

전체 경계 차단의 보조논증. \(T\) 를 식 (21)와 같이 \(H_{t}\)의 미소비
회계 증표와 미해결 이전 증표로 이루어진 정확한 미소비 경계
집합 전체에 대응하는 대상으로 두고, \(H_{r}≥H_{t}\)를 확정된 공통
확인 시점이라 하자. 식 (22)를 만족하는 조건은 다음과 같이 쓸
수 있다.

\begin{equation}
    T\ne\varnothing,\quad\forall t\in T:\operatorname{Blocked}_{H_r}(t).\tag{25}
\end{equation}
여기서 Blocked는 해당 시점에 미소비이면서 동결 또는 사용
금지 상태인 경우에만 참이다.

같은 나쁜 사건 밖에서 식 (25) 아래에서는 \(H_{t}\)와 \(H_{r}\) 사이에
시작점의 후속 경로를 새로 소비하는 전이가 수락되지 않았다.
또한 \(H_{r}\) 이후에도 \(T\) 의 차단 상태가 유지되는 동안 그 경로의
추가 소비는 수락되지 않는다.

첫 번째 결론에 반하는 전이가 있었다면, \(H_{t}\) 이후 최초의 그
러한 전이는 \(H_{t}\)의 미소비 경계 집합에 속한 한 객체를 소비해
야 한다. 그러면 그 객체의 불변 소비 기록은 \(H_{r}\)에 0일 수 없으
므로 식 (25)에 모순이다. 두 번째 결론에서도 최초의 추가 소
비는 차단된 미소비 경계 집합의 객체를 입력으로 필요로 한다.
모든 소비 경로가 현재 상태를 검사하므로 그 전이는 거부된다.
다른 정상 입력과의 병합이나 가공도 차단된 입력 하나의 검사
를 피할 수 없으며, 이미 종단인 제품 주장과 종료는 추가 소비
경로를 만들지 않는다.

완 료 출 력 게 임 에 서 는 공 격 자 가 합 법 적 전 이 · 상 태 변
경 의 스 케 줄 과 자 료 오 류 를 고 르 고 신 뢰 된 상 태 관 리 자
가 실 제 AuditAndFreeze를 실 행 한 다. 이 후 이 벤 트 까 지
다 항 예 산 에 포 함 하 며, \(H_{r}\)은 사 용 한 관 찰 · 영 수 증 을 반
영 하 는 같 은 확 정 이 력 의 공 통 시 점 이 다. Pr\([\)완료 반환 \(∧\)
    \((\)완료 술어 거짓 또는 차단 함의 위반\()]≤\epsilon _{B}\)이다. 나쁜 사건
밖에서는 추적 보조논증의 정확한 대상과 인증된 완료 검사에
위 귀납을 적용할 수 있기 때문이다. 이는 공동 사건의 비조건
부 상계이며 완료에 조건을 건 실패 확률 상계라고 주장하지 않
는다.

이 충분조건이 사용 차단의 모든 가능한 방식에 대한 필요조
건은 아니다. 예를 들어 \(T\) 의 객체가 종료로 소진되면 이 실행은
전체 완료를 보고할 수 없어도 그 객체에서 남은 소비 경로는 없
을 수 있다. 다른 대응 실행이 이미 후속 객체를 차단한 경우도
원래 \(T\) 의 완료와 구분해야 한다. 따라서 INCOMPLETE를 곧바
로 “현재 소비 가능한 후속 객체가 반드시 존재함”으로 해석하

% Source PDF page 14, column 1
지 않는다. \(T\) 가 빈 경우에도 프로토콜은 전체 완료와 구분하여
NO\_LIVE\_TARGETS를 반환한다.

정리 2의 통합 결론. 전이·추적·완료 논증은 같은 전체 실행의
나쁜 사건을 제외한다. 그 밖에서는 실제 생성 값에 대한 전이
가 유효하고, 추적의 성공 결과가 정확하며, 거짓 완료와 차단
유지 중 후속 소비가 모두 배제된다. 그러므로 통합 게임의 실
패도 이 하나의 나쁜 사건에 포함되어 식 (24)를 만족한다. 세
실패 상계를 독립 사건처럼 세 번 더하지 않는다. \(□\)

경합과 완료 판정의 반례. 서로 다른 시점의 성공 관찰. \(T=\)
\(\{a,b\}\)에 대해 먼저 \(a\)만 동결하고, \(a\)를 해제한 뒤 \(b\)만 동결했다
고 하자. 두 객체 각각에는 동결 성공 영수증이 있지만 둘이 동
시에 차단된 시점은 없을 수 있다. 따라서 과거의 객체별 성공을
단순 합산하여 현재 전체 완료를 보고할 수 없다. 공통 \(H_{r}\)에서
미소비와 상태를 확인하는 조건은 이 반례를 배제한다. 이는 해
제 가능한 동결을 포함한 본 절차의 조건이다. 모든 차단이 되
돌릴 수 없는 사용 금지이고 인증된 관찰을 이후 공통 원장 상태
로 정당하게 확장할 수 있다면 다른 확인 절차도 가능하므로, 모
든 프로토콜에 동일한 재조회가 필수라는 일반 명제는 아니다.

대상 축소와 소비 경합. \(a\)가 동결 전에 소비되어 새 미소비
객체 \(a^{′}\)를 만들었는데 \(a\)를 실행 결과에서 제거하고 \(b\)만 확인하

\(a^{′}\)
면,

의 차단 여부와 무관하게 거짓 전체 완료를 보고할 수 있
다. 본 절차는 \(T\) 를 고정하고 \(a\)의 소비 결과를 남기므로 이를 허
용하지 않는다. 유효한 소비가 동결보다 먼저 원장에 수락되면
해당 실행의 전체 완료는 성립하지 않으며, 새 기준 시점의 추적
으로 후속 상태를 조사해야 한다. 실행 자체는 유한 확인 한도
뒤 INCOMPLETE로 끝날 수 있다. 이는 전체 완료의 보장 부재
이지 모든 가능한 대응 방식의 종료 불가능성 정리가 아니다.

차단 사실과 운영상 완료. 식 (25)는 디지털 소비를 막는 충
분조건이다. AuditAndFreeze는 여기에 정확한 대상·결과 대
응과 불명확한 제출의 부재도 요구한다. 따라서 차단에 대한 상
태 성질과 실행 전체의 정산된 결과 보고는 구분된다. 이미 발
급된 제품 주장의 취소, 물리적 회수, 해제 이후의 계속된 차단
은 어느 조건에서도 따라오지 않는다.

감사 데이터 구성의 대안과 한계. 별도로 암호화한 부모 참조
와 출력의 미래 식별자를 모두 저장하는 것은 본 설계의 직접 탐
색 방식이다. 두 종류의 암호문이 모든 구성에서 기능적으로 필
수라고 주장하지 않는다. 예를 들어 관련된 모든 출력 기록과
과거 키 영역의 복호화 키를 확보하고, 공개 입력 식별자에 대응
하는 출력의 종류·위치를 식별할 수 있으면, 전체 출력의 미래
식별자를 복호화하여 다음 역색인을 만들 수 있다.

\begin{equation}
    M[k(r)]=r,\tag{26}
\end{equation}
여기서 \(k(r)\)는 종류를 포함한 객체 \(r\)의 소비 식별자이다. 입력
식별자를 이 색인에 대입하면 별도 부모 암호문 없이도 선행 객
체를 찾을 수 있다. 이는 식별자의 구별과 전체 관련 자료·키의
가용성에 조건부이며, 누락·중복 대응은 오류로 처리해야 한다.

반대로 전체 부모 참조로 소비 이벤트 역색인을 만드는 것만
으로는 아직 소비되지 않은 출력의 미래 식별자가 자동 주어지
지 않는다. 소유자 협력 없이 현재 객체의 상태를 변경하려면 그
집행 식별자를 확보하는 별도 수단이 필요하다. 이 비교는 직접
탐색과 전역 색인이라는 구성 선택을 설명할 뿐, 측정되지 않은
비용이나 효율 우월성을 뜻하지 않는다.

% Source PDF page 14, column 2

\subsection{프라이버시(Privacy)}
두 키 암호화의 기밀성 보조정리. 기밀성 이득은 두 실험에서
공격자가 1을 출력할 확률의 차이로 정의한다. 독립적으로 생성
한 A/B 공개키와 설정 전에 고정한 \(K−1\)개 이하의 위원회 조
각을 공개한다. 공격자는 같은 길이 감사 평문 쌍을 적응적으로
선택하고 한쪽의 정직한 암호문 및 \(W_{B}\)를 받는다. \(SK_{B}\), 복호
화·키 복구와 challenge wrapper의 부분 복호화 질의는 제공하
지 않는다. 인코딩 길이는 고정되고, 같은 길이의 공개 계산 가
능한 유효 dummy가 있다. 공개 descriptor와 인증은 각 환원
에서 재현 가능하며, 숨은 키 일치에 대한 별도의 비시뮬레이션
증거를 포함하지 않는다. 표준 다중 메시지 공개키 암호와 정적
조각 노출의 임계 암호 게임을 사용한다[16].

보조정리 (두 키 암호화). TE의 해당 CPA 이득을 \(\epsilon _{A}\), 데이터
PKE의 다중 메시지 CPA 이득을 \(\epsilon _{D}\)라 하면 이 게임의 이득은
\(2\epsilon _{A}+\epsilon _{D}\) 이하이다.

증명. 먼저 \(W_{B}\)를 실제 \(SK_{B}\) 대신 dummy의 암호화로 바
꾼다. TE 환원은 B 키쌍을 스스로 생성하므로 두 평문을 안다.
다음 데이터 암호문을 반대편 평문으로 바꾼다. PKE 환원은 받
은 \(PK_{B}\)의 \(SK_{B}\)를 몰라도 이미 dummy인 wrapper를 생성할
수 있다. 마지막으로 실제 키 wrapper를 복원한다. 세 변화의
삼각 부등식이 결론이다. 두 독립 키를 쓰므로 이 구조 자체에
자기 키 암호화 가정을 추가할 필요는 없다. 이 결과는 다른 공
개 필드의 비밀성이나 외부 키 전달 서비스의 인증·강건성을 증
명하지 않는다. \(□\)

원장 구별 게임과 충분조건. 전체 원장의 비밀성에는 커미트먼
트·소비 식별자·proof의 공동 분포도 필요하다. 특히 Voucher
는 같은 개방 난수를 commitment와 resolution secret에 사
용하므로 개별 hiding와 PRF의 이름만으로 그 합성을 결론 내
리지 않는다. 다음은 현재 수식에 대한 명시적 충분조건이다.
구체 해시·회로의 검증 완료를 뜻하지 않는다.

Note/Voucher commitment, resolution secret과 두 소비
식별자의 hash domain을 독립된 고전적 무작위 오라클로 모형
화한다. 출력 영역의 크기는 최소 \(2^{λ}\), 개방 난수와 정규 소유 비
밀의 독립 균등 영역은 최소 \(2^{κ}\)이고 \(κ\)는 \(λ\)에 선형으로 증가한
다. 소유 비밀 재사용은 허용하되 이 게임에서 소유 비밀·개방값
·별도 공개 OwnerKey는 제공하지 않는다. 공개 caller·서명키
등은 보호 비밀과 독립이다. 두 키 보조정리의 암호 가정과 함
께, 이 공동 오라클 모형에서 보조 입력에 안전한 다중 정리 ZK
및 모든 정규 statement에서 종료하는 PPT SimProof\((td,x)\)를
가정한다.

공격자는 공개 설정을 본 뒤 두 의미적 이력 \(h_{0},h_{1}\)을 고른
다. 이력은 회계 값·소유자 이름·연결을 지정하지만 비밀·난
수를 지정하지 않는다. 공개 누출 \(L\)이 같아야 한다. \(L\)은 종류
·정책· scope·타입·입출력 수·순서·sibling 관계·caller·시간·
기한·root의 prefix 위치·암호문 길이·공개 Claim의 문서 해시
와 논스 의미다. 공개 객체와 소비 식별자의 각 typed label,
producer/index 구조와 동일성 패턴도 포함한다. 비공개인 생
성 참조에서 미래 식별자로 가는 대응은 \(L\)에서 자동 연결하지
않는다. 원시 바이트가 같아야 한다는 요구는 아니다.

선택한 이력은 독립 정직 난수로 실현되며, 우연한 충돌 외에
는 그 난수에 따른 선택적 수락·비트 조건·재표본화를 요구하지
않는다. 관측자는 모든 공개 원장·암호문·증명을 보고 공개 오
라클 질의를 수행하지만, 참여자 부패·개방값 공개·능동 제출·
복호화 또는 보호 관계의 감사·상태 공개 질의는 받지 않는다.

% Source PDF page 15, column 1
위원회 노출은 설정 전 고정한 \(K−1\)개 이하로 유지한다.

정리 3 (Privacy). 위 게임에서 전체 proof의 ZK 이득을
\(\epsilon _{Z}\)라 하면

\begin{equation}
    \mathrm{Adv}_{\mathrm{hist}}\le2(\epsilon_Z+\epsilon_A+\epsilon_D+\epsilon_{\mathrm{RO}}).\tag{27}
\end{equation}
객체 수 \(n\), 소유자 수 \(u\), 외부·시뮬레이터의 관련 질의 수 \(Q\)에
대해 \(\epsilon _{RO}=O((n+u+Q)^{2}/2^{min(κ,λ)})\)의 보수적 상계를 사용
할 수 있다. 구체 보안 비트나 성능 수치가 아니다.

증명. 각 이력에서 먼저 참인 공개 문장의 proof를 ZK 시뮬레
이터로 바꾼다. 다음 두 키 보조정리의 순서로 wrapper와 감사
평문을 dummy로 치환한다. 마지막으로 숨은 commitment와
미래 식별자에 독립 무작위 label을 배정하되 공개 재사용·동일
성은 보존한다. Note commitment 입력에는 숨은 개방 난수
가, \(nf\) 입력에는 소유 비밀이 필요하다. 같은 비밀을 재사용해
도 다른 \(cm\)는 다른 오라클 입력이다. Voucher의 두 domain은
같은 개방 난수에 의존하나 resolution secret은 별도 균등 값
이다. 그 개방 난수 또는 resolution secret을 맞히기 전에는 공
동 오라클 표와 독립 label을 같은 분포로 결합할 수 있다.

비밀을 받지 않은 외부·시뮬레이터의 첫 숨은 입력 질의와 관
련 충돌을 나쁜 사건으로 두면 균등 비밀 추측 및 birthday 상계
가 \(\epsilon _{RO}\)를 준다. 그 밖의 root/index와 암호문 문맥은 공개 label
의 함수다. 마지막 분포는 동일 \(L\)과 독립 키·label로만 정해지
므로 양쪽 hybrid의 삼각 부등식으로 정리가 따른다. 암호문 치
환 후 statement가 거짓이어도 ZK를 다시 적용하지 않는다. 각
환원은 total PPT SimProof\((td,x)\)를 효율적 후처리로 실행할
뿐 중간 원장의 유효성을 주장하지 않는다. \(□\)

이 정리는 일반 참여자 프라이버시 전체와 다르다. 소유 비밀
을 공개하면 그 비밀로 과거·미래 Note 식별자를 대입할 수 있
고, Voucher opening을 아는 송·수신자는 그 해결 식별자를 안
다. 부패 참여자의 지식, 적응적 공개, 공개 OwnerKey를 허용
하는 확장과 실제 hash/증명계의 공동 모형 적합성에는 별도의
증명이 필요하다.
